\documentclass[11pt]{article}
\usepackage[utf8]{inputenc}
\usepackage{amsmath}
\usepackage{hyperref}

\title{Adaptive Real-Time Pump Scheduling}
\author{Gabriel Ribeiro}

\begin{document}
\maketitle

\section{Introduction}
Water distribution utilities face high electricity costs when operating pumping stations under time-of-use tariffs~\cite{bras2025}.
Day-ahead schedules optimize pumping costs, \textit{e.g.} by filling elevated reservoirs during off-peak hours.
As shown in \autoref{fig:arch}, the architecture separates economic planning from real-time feedback.

\begin{equation}
J = \sum_{k=0}^{N-1} c_k P_k \Delta t
\end{equation}

\section{Methodology}
It was developed a tracking controller that allows to correct pump speeds online.
The approach consists in a receding-horizon formulation that is different than conventional rule-based control..

\subsection{Experimental Setup}
We evaluate the WSS controller on a single-tank benchmark network with two variable-speed pumps.
And the WSS module updates setpoints every fifteen minutes while logging reservoir levels, hydraulic heads, energy consumption, and valve switching events across all simulated operating conditions under realistic stochastic demand disturbances.
For more details on the solver tolerances, \textit{etc.} see the appendix.


\section{Conclusions}
The proposed framework reduces energy cost without violating operational constraints.

\begin{thebibliography}{9}
\bibitem{bras2025} M. Bras et al., Smart-DLS optimization, 2025.
\end{thebibliography}

\end{document}
